\documentclass[11pt,a4paper]{article}
\usepackage[margin=2.35cm]{geometry}
\usepackage[T1]{fontenc}
\usepackage[utf8]{inputenc}
\usepackage{lmodern,microtype,booktabs,tabularx,enumitem,xcolor}
\usepackage[hidelinks]{hyperref}
\hypersetup{pdftitle={Community Dictionaries: Collaborative Resources and AI Maintenance},pdfauthor={ChatGPT C-LARA-Instance and Manny Rayner},pdfsubject={Revised stakeholder discussion draft, 26 September 2026}}
\setlist{nosep,leftmargin=*}
\setlength{\emergencystretch}{2em}
\newcommand{\repo}[2]{\href{https://github.com/mannyrayner/C-LARA-2/blob/e3b94b08d0d493e88e6966781323df2753f0bdfa/#1}{#2}}
\title{Community Dictionaries:\\Collaborative Language Resources\\and a Path to AI-Maintained Software}
\author{ChatGPT C-LARA-Instance and Manny Rayner\\C-LARA Project}
\date{26 September 2026\\\small Revised stakeholder discussion draft --- version 2}
\begin{document}
\maketitle
\begin{abstract}
Community Dictionaries is a lightweight web application for small groups to
construct and discuss picture and audio dictionaries. The initial version uses
phone photographs, human recordings and optional text, and has successful
human-reported trials on a laptop and a physical phone through the shared
C-LARA-2 AWS deployment. Following discussion with Sophie Rendina, we propose
optional AI image generation and text-to-speech as the next increment, allowing
comparison with human-produced media. Community authority over data is a design
requirement for these options. A second, longer-term aim is to reduce dependence
on a technical intermediary: linguists should be able to request changes
directly, while an AI system handles increasingly complete maintenance tasks.
We outline a staged experiment in reducing Manny's operational involvement
from daily intervention to several reviews per week and then per month.
The implemented prototype, proposed extensions and untested autonomy goals
are distinguished throughout.
\end{abstract}

\section{A working starting point}
The application's purpose is to help people build useful language resources
together with little technical overhead. Indigenous and other under-resourced
languages are a central motivation, but the same design can support languages
with mature AI services. Swedish provides an initial language whose content
Manny can assess confidently; an Icelandic pilot with Kate and Axel has been
invited. Neither establishes suitability for an Indigenous community.

Construction and discussion come first. An entry can begin with a photograph
and acquire a recording from another member. Optional fields hold a written
form, meaning or translation, and category. Members can browse pictures, search
supplied text, listen, and add written or recorded comments. Named partnerships
organise requests for missing pictures or recordings, making shared work easy
to retrieve. Editors review contributions; accepted wording is versioned and
alternative media retain their attribution.

Each dictionary has an invited membership. Partnerships organise work within
that dictionary; a private subgroup needs a separate dictionary. Media is
stored outside public web directories and served through membership checks.
These controls do not exclude technical access by server operators. An export
contains records and media; operational recovery also requires the shared
accounts, database and private-file backups.

The app is available at
\url{https://c-lara-2.c-lara.org/community-dictionaries/}.
It is a separate Django application sharing C-LARA-2's accounts and deployment.
It uses a small browser interface, direct camera/microphone input and uploaded
files, without invoking the text-compilation pipeline. The current version has
\textbf{no AI content generation}. Optional AI media described below is the
next proposed increment. Dedicated practice exercises and short video clips
remain later work.

\section{Development and evidence so far}
The assistant wrote the initial implementation, tests and documentation in two
substantial iterations; Manny supplied goals, scope decisions, installation
and feedback. No separate coding-agent handoff was used. Each iteration
included multiple tool calls, tests and repairs. The development conversation
identifies the assistant as GPT-6 Astra, but the iteration records do not
independently capture the serving model identifier or token cost.

Manny estimated the first functional cut at 45 minutes and, in his subsequent
email, the repair at another 20 minutes. These are participant estimates, not a
measured total including specification, existing infrastructure, installation,
evaluation and maintenance. Human feedback mattered: initially hidden text
fields were made prominent, and microphone selection and signal feedback were
improved. Manny resolved his silent recording by choosing the correct input
device. He later made the one-line production change disabling Django debug.

\begin{table}[htbp]
\centering\small
\renewcommand{\arraystretch}{1.12}
\begin{tabularx}{\linewidth}{@{}p{0.23\linewidth}X@{}}
\toprule
Evidence & Scope of the result \\
\midrule
App tests & Iteration 2 records 21 passing Django tests covering workflows and access, including retries, review, private media and export/fixture restoration. This is not a production restore rehearsal. \\
Browser rehearsal & Controlled Chromium checks exercised capture, retry, search, review and silent-retake handling. Mobile-sized viewports were not physical phones or Safari. \\
Human use & Manny reported successful AWS laptop use on 25 September. Cathy then took and saved a photo on her phone; Manny added a recording. The phone/browser and the device used for the recording were not recorded. \\
Repository Assistant & On 26 September Manny confirmed that, after the documentation update, the server Assistant correctly distinguished this phone success from outstanding trials and deployment work. \\
\bottomrule
\end{tabularx}
\caption{Available evidence. Previous automated tests were not rerun to prepare this report.}
\end{table}

These results support an initial functioning system. They do not establish a
complete device matrix, sustained adoption, learning gains or long-term
maintenance reliability. Production work also remains: \texttt{DEBUG=False}
is confirmed, but secure-cookie settings, host restrictions and HTTPS/HSTS
policy need completion. An external check confirmed HTTP-to-HTTPS redirection
for the dictionary address; Django's deployment checker cannot inspect nginx's
complete policy. A successful functional trial is not a security review.

The Assistant's corrected answer demonstrates a useful maintenance mechanism:
recording events in the repository lets a later conversation recover them.
Its current execution mode is explicitly read-only. Explaining the project
accurately and operating it autonomously are different capabilities.

\section{Next increment: optional generated media}
Sophie proposed testing current image generation on meanings that are difficult
to illustrate. The next version should allow comparison of photographs and
generated pictures, alongside comparison of human recordings and TTS where
suitable voices exist. A preference for photographs would itself be a useful
finding. Generation should be optional per dictionary and per operation, with
camera and human recording remaining fully usable when AI is disabled.

\paragraph{Images.}
A participant selects an entry, describes the intended meaning and context,
checks the material to be sent, and requests a candidate image. Guidance can
be supplied in an explanation language when the model does not understand the
target language. Only that authorised prompt should leave the service; existing
photos, recordings and unrelated entries should not be attached implicitly.
The result is previewed and submitted for review with its AI origin recorded.
It must not silently replace an accepted photograph.

\paragraph{Speech.}
A participant supplies or confirms text, selects an approved language/voice,
and generates a recording for listening and review. Language availability
must be distinguished from acceptable pronunciation, dialect and intended
meaning. Swedish is a practical first check; Icelandic requires an appropriate
speaker's assessment, and neither should imply support for an Indigenous
language. Unsupported or failed synthesis must leave human recording available.
There is no proposed voice-cloning feature.

\paragraph{Reuse and remaining engineering.}
Source inspection confirms an image-generation client in
\repo{src/core/ai_api.py}{\texttt{src/core/ai\_api.py}} and OpenAI/Google TTS
engines in \repo{src/pipeline/audio.py}{\texttt{src/pipeline/audio.py}}.
We can adapt these without running the full compilation pipeline. However,
new integration is needed for dictionary permissions, private output storage,
provider policy, job status, cost limits and provenance. The audio pipeline
also has a test-tone fallback; a user-facing dictionary adapter must require
real synthesis and report failure explicitly.

Generated contributions should record the requesting member, source text or
prompt, provider/model, voice where relevant, and review history. Long-running
requests need bounded jobs so phone interactions remain responsive. A retry
must not create duplicate contributions or automatically repeat an uncertain
paid provider request. Existing user-key and billing mechanisms offer reusable
parts, but community-approved provider restrictions must take precedence over
a user's preferred API key. No live provider calls were made for this design
review, and neither new media feature is yet deployed.

\section{Data sovereignty as a design requirement}
Sophie's response places data sovereignty at the centre of viability. Our
control of the hosting and code gives us useful implementation choices; it
does not establish community authority by itself. The Maiam nayri Wingara
principles describe Indigenous control extending across the data lifecycle
and its infrastructure. The CARE principles likewise foreground collective
benefit, authority, responsibility and ethics \cite{mnw,care}. These are useful
starting points for discussion with the particular community, not a substitute
for that discussion.

For this app, the proposed practical outcome is a short, community-approved
policy identifying who decides membership, publication, export, deletion,
external processing and acceptable reuse. It should cover photographs and
voices as well as language text, and account for restricted material, backup
retention, browser drafts and copies already downloaded. A dictionary owner
account is an application role; it is not proof of authority to decide on
behalf of a community. Separate hosting or a community-controlled deployment
may be appropriate where the shared service cannot meet requirements.

The first AI integration should keep external services off unless permitted
for that dictionary and explicitly invoked. Where zero data retention or other
provider conditions are required, the actual account, endpoint and configuration
must be verified; an unverified service remains unavailable. A no-training
condition and a no-retention condition are distinct. An English explanation can
also reveal restricted knowledge, so translation alone does not remove the need
for authorisation.

There are two external-data paths to govern: generating dictionary content and
using an AI to diagnose or maintain the software. Disabling image generation
must not be undermined by sending private entries or raw production logs to a
maintenance model. Routine diagnosis should use code, synthetic examples and
minimal redacted diagnostics. These controls need implementation and testing;
they are not claims about a completed governance system.

\section{Reducing dependence on a technical intermediary}
Manny's central longer-term proposal is to remove himself from routine technical
operations. Today he installs patches, checks traces, deploys changes and
translates collaborators' requests into development work, often daily. The
proposed progression is several contacts per week, then several per month,
and subsequently a review of what further autonomy is realistic. Cathy supports
the objective while questioning whether it can be achieved. That makes a
measured experiment more useful than a promise about a future model.

The desired workflow is concrete: a linguist reports a problem in ordinary
language; the AI clarifies the intended behaviour with that person, reproduces
the problem, prepares a change and meaningful checks, presents a preview, and
arranges release under an agreed operating policy. It then monitors the result,
recovers from an unsuccessful release where possible, and records what changed.
The linguist judges whether the result meets the need. Manny need not relay
each exchange or execute each command.

This requires an operational service around the model: persistent work queues,
repository access, an isolated build/test environment, staging, a release
mechanism, monitoring, budget controls, backup and restore procedures, and
credential renewal. The AI needs evidence of actual deployments and failures,
not just a memory of suggested commands. The current read-only Assistant could
remain the explanation interface while a separately authorised execution
service performs defined work. Its permissions should be enforced by the
service; natural-language requests cannot confer additional authority.

\begin{table}[htbp]
\centering\small
\renewcommand{\arraystretch}{1.13}
\begin{tabularx}{\linewidth}{@{}p{0.19\linewidth}p{0.24\linewidth}X@{}}
\toprule
Stage & Manny's routine role & Evidence needed before reducing involvement \\
\midrule
Baseline & Current daily intervention & Log technical contacts, minutes, blockers, incidents and total human effort for two weeks, alongside actual workload. \\
Batched supervision & Two or three scheduled reviews per week & Collaborator requests reach the AI directly; ordinary fixes arrive reproduced, tested and previewable. Demonstrate release/rollback and restoration in a safe test environment. \\
Routine delegation & Two scheduled reviews per month & Previously authorised, reversible change classes can be released and monitored without manual patch installation. Observe a representative month plus recovery exercises. \\
Exception-led operation & Periodic strategic review; exceptional escalation & Routine operations continue during Manny's planned absence; failures reach a named fallback when automated recovery cannot safely finish. \\
\bottomrule
\end{tabularx}
\caption{Proposed stages, not delivery dates or achieved service levels. Advancement depends on results; a failed stage should be revised or repeated.}
\end{table}

The measures should include interruption count, hands-on minutes, resolution
time, unresolved backlog, service incidents, failed releases and operating
cost. Fewer messages alone is not success if requests are ignored, work is
batched into exhausting sessions, or Sophie and Cathy inherit the technical
burden. Comparable request volume and complexity matter. Quiet weeks should
be supplemented with controlled exercises such as failed deployment, restore,
expired credentials and provider unavailability, without disrupting live users.

Routine reversible operations can eventually run under standing authorisation.
Changes to community data policy, access boundaries, spending limits or destructive
migrations need an authorised decision or a previously approved procedure.
That decision need not always belong to Manny. Communities retain authority
over their resources; an accountable organisation still needs to own hosting,
provider and billing arrangements. A named emergency contact is necessary for
failures beyond the system's recovery abilities. The goal is to remove dependence
on one technical person without silently transferring it to another.

The current deployment has no autonomous maintenance service. A stronger model
may improve diagnosis and implementation, but reliable
operation also depends on the surrounding services and tested recovery paths.
A bounded reduction in manual work can be tested with present tools; the claim
of fully AI-maintained operation remains an open empirical question.

\section{Next discussion and evaluation}
A discussion with Sophie is planned for early afternoon on Thursday 1 October.
It should select a small set of useful meanings, agree which material may be
used with external services, and identify who can approve the relevant policies.
A first comparison might cover familiar objects, actions and harder-to-depict
meanings, using photographs and AI candidates for the same senses. Review should
address recognisability, cultural/contextual fit, ambiguity, correction effort
and time to an acceptable entry, rather than visual attractiveness alone.
TTS should similarly be compared with human recordings for intelligibility
and pronunciation. These are proposed qualitative trials, not reported results.

The Icelandic trial can examine collaboration and the desire to return to the
app. Indigenous-community testing requires its own participants, priorities
and governance decisions. For usability, the ambition is a contribution flow
comfortable for people with limited literacy; optional text is only a starting
point, and spoken explanations, icons and reduced navigation may need further
work with users. Maintenance evaluation should run alongside feature work,
starting with an intervention log and one complete request-to-release example.

\section*{Contributions and evidence}
ChatGPT C-LARA-Instance prepared the software and report drafts. Manny provided
project direction, installation, evaluation reports and the debug-setting
change; Cathy contributed the phone trial. Sophie's correspondence motivated
the sovereignty emphasis and image-comparison proposal. This draft does not
imply approval by a community or endorsement by future testers.

The reviewed source is commit
\href{https://github.com/mannyrayner/C-LARA-2/commit/e3b94b08d0d493e88e6966781323df2753f0bdfa}{\texttt{e3b94b0}}.
The repository contains the
\repo{experiments/community_dictionary/iteration-001.md}{first iteration},
\repo{experiments/community_dictionary/iteration-002.md}{repair iteration}, and
\repo{experiments/community_dictionary/aws-phone-trial-2026-09-25.md}{AWS/phone trial}.
The accompanying roadmap specifies reuse boundaries, acceptance checks and the
proposed autonomy experiment. This second report pass changes documentation;
AI media and an autonomous execution service remain to be implemented.

\begin{thebibliography}{9}
\bibitem{mnw} Maiam nayri Wingara. \emph{MnW Principles}: principles from the
2018 Indigenous Data Sovereignty Summit. Accessed 26 September 2026.
\url{https://www.maiamnayriwingara.org/mnw-principles}.
\bibitem{care} Global Indigenous Data Alliance. \emph{CARE Principles for
Indigenous Data Governance}. Accessed 26 September 2026.
\url{https://www.gida-global.org/careprinciples}.
\end{thebibliography}
\end{document}
